%==========================================================
\section{Organização da equipa}\label{anx:organizacao}

Cada tipo de modelo tem um dono, que define a notação final e revê todos os modelos desse tipo, e um suplente, para que o conhecimento não dependa de uma só pessoa.

\begin{table}[H]
\centering\small
\begin{tabularx}{\textwidth}{@{}p{3.6cm}Lp{2.2cm}p{2.2cm}@{}}
\toprule
\textbf{Responsabilidade} & \textbf{Âmbito} & \textbf{Dono} & \textbf{Suplente} \\
\midrule
Fluxo & Notação, exemplo e revisão de todos os fluxos & Tiago & Gonçalo \\
Cultural & Idem, e das «vozes» das setas & Tiago & Gonçalo \\
Sequência & Idem, e das intenções e gatilhos & Gonçalo & Zé Pedro \\
Artefacto & Idem, e do arquivo anonimizado de artefactos & Gonçalo & Zé Pedro \\
Físico & Idem & Zé Pedro & Tiago \\
Guardião do processo & Calendário, rotação de papéis, consentimentos, repositório & Zé Pedro & Tiago \\
\bottomrule
\end{tabularx}
\caption{Donos e suplentes de cada modelo.}
\label{tab:donos}
\end{table}

Nas sessões de interpretação os papéis rodam, para que todos desenhem todos os tipos de modelo e ninguém modele sozinho a sua própria entrevista. Só o modelo físico é desenhado pelo entrevistador.

\begin{table}[H]
\centering\small
\begin{tabularx}{\textwidth}{@{}p{3.6cm}Lp{2.2cm}p{2.2cm}@{}}
\toprule
\textbf{Papel na sessão} & \textbf{Desenha} & \textbf{A0} & \textbf{U0} \\
\midrule
Entrevistador(a) & Físico; conta a entrevista pela ordem & Tiago & Gonçalo \\
Modelador(a) A & Fluxo e cultural; acrescenta notas & Gonçalo & Zé Pedro \\
Modelador(a) B & Sequência e artefacto & Zé Pedro & Tiago \\
\bottomrule
\end{tabularx}
\caption{Rotação de papéis nas sessões de interpretação.}
\label{tab:rotacao}
\end{table}

\pendente{Equipa}{Confirmar quem entrevistou A0 e U0 e ajustar a Tabela~\ref{tab:rotacao}. Acrescentar a contribuição de cada elemento para esta entrega (o enunciado prevê um inquérito de contribuição dentro da equipa).}
